\section{LASA 微架构}
\label{sec:microarch}

\subsection{沿数据流组织的硬件结构}

\cref{fig:datapath} 按一个 tile 的数据流给出 \LASA 结构。bias、weight 和
IFM 由三个独立 64 bit AXI-Stream 输入，OFM 通过第四路 DMA 写回。bias 和
weight loader 检查 TKEEP、TLAST 和声明字节数，并将每个 32 通道块的 bias
及当前/下一上下文权重写入本地缓冲。IFM loader 先执行
\begin{equation}
x_{s8}=\mathrm{sat}_{[-128,127]}(x_{u8}-z_x),
\end{equation}
再进入原生 1$\times$1 向量路径或 3$\times$3 HWC 窗口物化路径。数组结果
由 collector 形成 32-lane int32 packet：非末 pass 写入 partial-PSUM 环，
末 pass 依次经过重量化、激活 LUT、可选 pooling 和 packed-HWC 输出。

\begin{figure*}[t]
\centering
\resizebox{0.98\textwidth}{!}{%
\begin{tikzpicture}[>=Latex,node distance=5mm and 5mm,
  b/.style={draw,rounded corners,minimum height=9mm,minimum width=20mm,align=center,font=\small},
  f/.style={->,thick}, fb/.style={->,thick,LasaRed,dashed}]
\node[b,fill=LasaGreen!15] (ifm) {IFM AXIS\\中心化};
\node[b,right=of ifm,fill=LasaCyan!15] (mat) {HWC窗口\\物化/银行化};
\node[b,right=of mat,fill=LasaCyan!15] (replay) {tile cache\\重放};
\node[b,right=of replay,fill=LasaBlue!18] (array) {18$\times$16\\双通道阵列};
\node[b,right=of array,fill=LasaOrange!18] (drain) {PSUM\\collector/drain};
\node[b,right=of drain,fill=LasaOrange!18] (rq) {requant\\LUT激活};
\node[b,right=of rq,fill=LasaGreen!15] (pack) {packed HWC\\OFM AXIS};
\node[b,below=10mm of array,fill=gray!15] (weight) {双上下文\\权重预载};
\node[b,below=10mm of drain,fill=LasaRed!10] (psum) {continuous PSUM\\ping-pong/feeder};
\draw[f] (ifm)--(mat); \draw[f] (mat)--(replay); \draw[f] (replay)--(array);
\draw[f] (array)--(drain); \draw[f] (drain)--(rq); \draw[f] (rq)--(pack);
\draw[f] (weight)--(array); \draw[fb] (drain)--(psum); \draw[fb] (psum)--(array);
\draw[decorate,decoration={brace,amplitude=5pt,mirror},thick]
  ($(mat.south west)+(0,-17mm)$)--($(replay.south east)+(0,-17mm)$)
  node[midway,below=6pt,font=\small]{materialize once, replay many};
\end{tikzpicture}
}
\caption{\LASA 的主数据通路。蓝色链路执行当前 pass，红色虚线为非末 pass
的 partial-PSUM 反馈；权重预载和 HWC 重放与当前计算并行推进。}
\label{fig:datapath}
\end{figure*}


结构上的关键不是模块数量，而是数据路径与控制路径的隔离。像素、权重和
PSUM 在规则流水中只依赖 ready/valid；context descriptor、epoch、tile 坐标、
pass/Cout 索引和错误状态在较窄控制网络中传播。r3 的路由分析曾显示，一个
descriptor compare 到 2048 bit payload clock-enable 的高扇出组合锥占据大量
负裕量。r4/r5 将 drain 改为静止双槽环：宽 payload 只在 FIFO return 时写，
ready 仅更新窄指针和计数。类似地，materializer 在上下文入口锁存静态准入
摘要，而动态读冲突仍逐拍检查。该经验说明，面向 200 MHz 的微架构必须避免
让协议判定直接驱动宽数据寄存器使能。

\subsection{HWC 窗口物化与银行映射}

3$\times$3 路径直接接收 HWC 字节。对输出位置 $(o_y,o_x)$、当前 pass 起点
$k_b$ 和 lane $r$，全局 $k=k_b+r$ 被拆为
\begin{equation}
\begin{aligned}
c_i&=\lfloor k/9\rfloor, & q&=k\bmod 9, \\
k_y&=\lfloor q/3\rfloor, & k_x&=q\bmod3.
\end{aligned}
\end{equation}
由此得到输入坐标 $(i_y,i_x)$。越界位置输出内部 signed zero；有效位置由
row store 和 byte bank 取数。由于通道在 HWC 中连续，bank 映射把相邻字节
散到独立端口，并以 row/tag metadata 判断所需行是否属于当前 epoch。
\cref{fig:hwc-banks} 展示缩略映射。

\begin{figure*}[t]
\centering
\begin{tikzpicture}[>=Latex,node distance=7mm,
  byte/.style={draw,minimum width=9mm,minimum height=7mm,font=\scriptsize},
  bank/.style={draw,rounded corners,minimum width=28mm,minimum height=10mm,align=center,fill=LasaCyan!12}]
\foreach \i/\lab in {0/{p0,c0},1/{p0,c1},2/{p0,c2},3/{p0,c3},4/{p1,c0},5/{p1,c1},6/{p1,c2},7/{p1,c3}}{
  \node[byte,fill=LasaGreen!10] (x\i) at (1.0*\i,0) {\lab};
}
\node[anchor=east] at (-0.3,0) {DDR HWC字节流};
\node[bank] (b0) at (1.1,-1.7) {Bank 0\\$(pixel,channel)\bmod B=0$};
\node[bank] (b1) at (4.3,-1.7) {Bank 1\\$(pixel,channel)\bmod B=1$};
\node[bank] (meta) at (7.5,-1.7) {row/tag metadata\\valid/epoch/地址};
\foreach \i in {0,...,7}{
  \pgfmathtruncatemacro{\target}{mod(\i,2)}
  \ifnum\target=0 \draw[->,LasaBlue] (x\i.south)--(b0.north);\else \draw[->,LasaBlue] (x\i.south)--(b1.north);\fi
}
\node[draw,rounded corners,fill=LasaOrange!12,minimum width=84mm,minimum height=11mm,align=center] (win) at (4.3,-3.5) {窗口地址生成：$(o_y,o_x,k_{base},lane)\rightarrow(row,bank,byte)$\\越界直接注入内部 signed zero};
\draw[->,thick] (b0)--(win); \draw[->,thick] (b1)--(win); \draw[->,thick] (meta)--(win);
\node[draw,rounded corners,fill=LasaBlue!15,minimum width=55mm,minimum height=10mm,align=center] (vec) at (4.3,-5.0) {每个输出像素的 18-lane $K$ 向量\\写入 replay cache，跨 pass 重放};
\draw[->,thick] (win)--(vec);
\end{tikzpicture}
\caption{HWC 输入的银行化、窗口物化和重放。硬件在片上保持 HWC 语义，
不要求 PS 生成完整 im2col；地址和 epoch 元数据用于防止跨上下文覆盖。}
\label{fig:hwc-banks}
\end{figure*}


物化器不会生成完整 im2col 矩阵，而是只为当前 spatial tile 生成
$S_T\times N_K$ 个 18-lane 向量。每个向量带像素顺序和 pass 身份写入
materialized cache。后续 Cout block 可以重用同一组输入；同一 Cout block
中的后续 pass 也从 cache 顺序重放。缓存深度门禁为
\begin{equation}
S_T N_K\le D_M=32768.
\end{equation}
对 m19，$S_T=26\times6=156$、$N_K=192$，乘积 29952；tile 高度 7 将超过
门限。缓存采用 URAM 承担长 tile 数据，BRAM/LUTRAM 负责浅 FIFO、metadata、
packet 和小型重排。这样的职责分离避免用 LUT 构造大容量位存储，同时让
context valid/tag 检查靠近小型控制存储。

物化与重放之间存在三条安全不变量：尚未完成的 reader 所使用的行不得被新
writer 覆盖；held 上下文的 active read rows 和 spec tag 每周期重新判断，
不能只在准入时冻结；提交 cursor 即使遇到越界 lane 也必须按描述符推进，
以保持后续地址一致。这些不变量在随机 hold、defer、row-start、重叠上下文
和边界 lane 测试中由软件 oracle 逐拍核对。

\subsection{18$\times$16 双输出通道阵列}

实际阵列由 18 行、16 个物理列组成。每个物理列的 DSP48E2 链同时维护两个
输出通道的 weight 和 PSUM lane，因此每个 pass 覆盖 32 个输出通道。
对阵列行 $r$ 和物理列 $j$，PE 执行
\begin{align}
p^{(0)}_{r,j}&=p^{(0)}_{r-1,j}+x_{r,j}w^{(0)}_{r,j},\\
p^{(1)}_{r,j}&=p^{(1)}_{r-1,j}+x_{r,j}w^{(1)}_{r,j}.
\end{align}
IFM 在水平方向传播，PSUM 沿 DSP cascade 纵向传播，权重在上下文期间驻留。
每拍最多完成 $18\times16\times2=576$ 次 int8 MAC。200 MHz 下理论峰值为
115.2 GMAC/s；若按一次乘和一次加计两次运算，则为 230.4 GOPS。

\begin{figure*}[t]
\centering
\begin{tikzpicture}[>=Latex,scale=0.9,transform shape]
\foreach \r in {0,...,4}{
  \foreach \c in {0,...,5}{
    \node[draw,minimum width=9mm,minimum height=7mm,fill=LasaBlue!10] (p\r\c) at (1.1*\c,-0.85*\r) {PE};
  }
}
\foreach \r in {0,...,4}{
  \draw[->,LasaGreen!70!black,thick] (-1,-0.85*\r)--(p\r0.west);
  \foreach \c in {0,...,4}{\draw[->,LasaGreen!70!black] (p\r\c.east)--(p\r\the\numexpr\c+1\relax.west);}
}
\foreach \c in {0,...,5}{
  \draw[->,LasaOrange!90!black,thick] (p0\c.north)+(0,0.6)--(p0\c.north);
  \foreach \r in {0,...,3}{\draw[->,LasaOrange!90!black] (p\r\c.south)--(p\the\numexpr\r+1\relax\c.north);}
  \draw[->,LasaRed,thick] (p4\c.south)--+(0,-0.6);
}
\node[align=left,anchor=west] at (7.4,-0.4) {实际阵列：18 行$\times$16 列\\每列两套 weight/PSUM lane\\每拍最多 576 次 INT8 MAC};
\node[draw,rounded corners,fill=LasaLight,align=left,anchor=west] at (7.4,-2.3) {
PE：$p'_{0}=p_{0}+xw_{0}$\\\phantom{PE：}$p'_{1}=p_{1}+xw_{1}$\\
权重驻留；IFM 横向；PSUM 纵向};
\node[LasaGreen!70!black] at (-1.2,0.55) {IFM};
\node[LasaOrange!90!black] at (2.8,0.8) {bias/feedback PSUM};
\node[LasaRed] at (2.8,-4.2) {32-lane PSUM packet};
\end{tikzpicture}
\caption{18$\times$16 双输出通道脉动阵列与 PE 语义。图中缩略为 5$\times$6，
实际每个物理列并行计算两个输出通道，因此 16 列形成 32 通道输出块。}
\label{fig:array-pe}
\end{figure*}


权重装载采用双上下文存储。当前 bank 计算时，下一 bank 可以接收下一个
pass 或下一个 Cout block 的权重；只有当所有权 scoreboard 证明当前 bank
已退休、下一描述符完整且 context tag 匹配时才允许切换。DSP cascade 使
PSUM 沿相邻 DSP48E2 的专用链传递，减少 LUT 加法树和通用互连。正式实现
中阵列占 576 个 DSP，全系统为 650 个 DSP；二者必须在资源表中区分。

18 行与 16 列并非只由器件 DSP 总量决定。行数决定单个 pass 覆盖的 $K$
元素和 DSP cascade 深度；列数决定每拍并行的输出通道块。继续增加行数能
直接降低 m13/m19 的 pass 数，但会延长输入广播与纵向级联路径，并增加每个
上下文的权重宽度；增加列数则扩大 PSUM packet、requant lane 和写回交叉连接，
对已经接近拥塞上限的 CLB 区域不友好。双输出通道让一个物理列承担两个逻辑
通道，在不把 IFM 广播网络翻倍的情况下得到 32-lane 输出，适合以 32 通道
为单位的量化参数、bias 和 HWC 打包。其代价是尾块必须独立屏蔽，且所有权
和错误计数都要以逻辑 32 lane 而非物理 16 列定义。

阵列采用完全确定的 token 顺序，而不是允许 PE 自主跳过零值。这样做牺牲了
非结构化稀疏性的潜在收益，却简化了跨 pass 对齐：给定 context tag 和 pixel
序号，反馈 PSUM、输入向量与权重向量在每一行都有唯一配对。确定性顺序也是
byte-exact 验证的基础；若硬件根据数据动态改变执行次序，功能结果可能仍然
相同，但性能计数、异常恢复和复现实验会变得依赖输入分布。本文只为结构上
确定的 1$\times$1 路径消除无用 3$\times$3 位置，不声称支持动态稀疏。

\subsection{原生 1$\times$1、3$\times$3 与尾 lane}

若把 1$\times$1 卷积机械地映射为中心位置非零的稀疏 3$\times$3，会将
$K=C_{in}$ 扩大为 $9C_{in}$，同时浪费输入窗口和权重传输。\LASA 提供原生
1$\times$1 vector path：每个像素按 18 个通道组成 pass，直接进入 IFM FIFO，
不经过 3 行窗口。m14、m16 和两个检测头使用该路径，其中 P5 的
$C_{in}=512$ 需要 29 个 pass，而稀疏 3$\times$3 映射需要 256 个 pass。

对于 $K$ 尾 pass，超过 $K_{total}$ 的 lane 注入零权重和零 IFM，valid/tag
仍按完整 18 行流水推进。对于 Cout 尾块，lane 有效位由
$c_{base}+lane<C_{out}$ 决定；无效 lane 不产生输出字节。P4/P5 的 255 通道
由 8 个 32 通道块执行，最后一块 31 lane 有效。packed-HWC writer 维护跨
packet 的 byte reservoir，只有全帧最后 beat 使用相应 TKEEP/TLAST，因而
每个像素 255 B 的非 8 字节对齐不会插入 padding。

原生 1$\times$1 与尾 lane 共同体现“层自适应”不是简单选择 tile 高度。
1$\times$1 路径改变的是 $K$ 轴生成方式，尾 pass 处理的是行方向不整除，
尾 Cout 块处理的是输出通道方向不整除，而 byte reservoir 处理的是 64-bit
总线方向不整除。四种边界必须组合验证：例如 P4 检测头既有 255 通道尾块，
又跨两个空间 tile，最后一个像素的最后 7 B 才触发最终 TKEEP。若仅对每个
32-lane packet 做局部对齐，就会在像素之间插入填充，主机仍可能读到正确
总字节数，却会从第一个像素之后全部错位。因此 writer 的完成条件由全帧
有效字节计数而非 packet 数推导。

\subsection{PSUM 反馈、重量化和写回}

第一 pass 的阵列顶部注入 int32 bias；非末 pass 的 32-lane 结果按
$(pixel,c_{block})$ 顺序进入 PSUM ping-pong。下一 pass 的 feeder 在阵列可
接受新 token 时读取相同顺序的 packet。所有权 scoreboard 跟踪分配、读、
写和释放，禁止同一 bank 同时被两个 context 占有。drain ring 使用两个静止
payload 槽，return/pop 的四种组合只更新窄指针和 count；满且同拍 pop 时
可以预留下一拍 return，从而维持稳态一 packet/clock，同时消除 ready 对宽
payload CE 的高扇出。

末 pass 的 int32 结果使用层级 per-tensor 参数重量化：
\begin{equation}
q=\mathrm{sat}_{s8}\left(\left\lfloor
\frac{p\cdot m+2^{14+s}}{2^{15+s}}\right\rfloor+z_y\right),
\end{equation}
其中 $m$ 为 16 bit multiplier，$s$ 为移位。为与 RTL 完全一致，host golden
复现正 half 加法、算术右移、signed clamp 和 256 项 LUT；它不直接采用框架
的 round-to-even。输出 byte 经可选 activation/pool 后按
\begin{equation}
a=((p_{global}C_{out})+c_{base}+lane)
\end{equation}
写回 HWC 线性地址。packet FIFO、byte reservoir 与 AXIS backpressure 共同
保证地址连续、尾字节压紧且 TLAST 只在预期总字节数处出现。

\subsection{背压与错误封闭}

四路 AXIS 都可能独立反压。数据面只有在 ready\&valid 时推进，context
cursor 只有在对应资源真实提交后更新。输入长度、TKEEP/TLAST、描述符容量、
时钟、内部 endpoint、DMA 状态和输出字节数均有 fail-close 检查；任一 sticky
error 会阻止该层被记为成功。形式化实现门禁还要求所有 routable net 完全
路由、setup/hold/pulse endpoint 为零、unconstrained/unclocked endpoint 为
零、DRC Error/Critical Warning 为零。这样，软件看到的 done 不只是阵列停止，
而是数据、协议和物理实现都满足发布合同。

片上存储按访问模式分工。URAM 保存深而规则的 materialized 向量和长层 tile，
其容量效率高但端口与物理位置固定；BRAM 保存权重、PSUM packet、行缓冲和
中等深度 FIFO；LUTRAM 用于短队列、索引和高端口小表。把所有缓冲统一映射到
一种存储会简化参数化，却会产生明显的物理代价：用 LUT 构造宽深 payload
增加控制扇出，用 URAM 放小 FIFO 浪费容量且限制布局。本文的划分原则是让
宽 payload 的写使能尽量来自局部注册状态，让跨模块 descriptor 只驱动窄
tag、pointer 和 count。r4/r5 时序收敛过程表明，消除宽寄存器 clock-enable
上的组合 ready 链，比继续缩短已经流水化的数据 D 路径更有效。

背压处理同样遵循“状态提交晚于数据所有权确认”。AXIS valid 可以保持，
但地址 cursor、tile 完成和 context retirement 只在实际 handshake 后推进；
错误路径不能借用正常 done 清理尚未提交的数据。reset、idle、start、正常
return/pop 与 final completion 具有固定优先级，异常同拍到达的数据即使写入
物理槽，也会被清零的 occupancy 遮蔽而不可见。这些规则看似增加控制逻辑，
却允许软件在收到错误后确定当前输出不可使用，而不是面对“部分更新、状态却
显示完成”的模糊故障。
